\chapter{離婁章句上}
\kaishi*

\jiazhu[1]{離婁者，古之明目者，蓋以爲黃帝時人也。黃帝亡其玄珠，使離朱索之。離朱即離婁也，能視於百步之外，見秋毫之末。然必須規矩，乃成方員。猶《論語》「述而不作，信而好古」，故以題篇。}

\section{論規矩仁政}

孟子曰：「離婁之明，公輸子之巧，不以規矩，不能成方員。」\jiazhu[1]{公輸子，魯班，魯之巧人也。或以爲魯昭公之子。雖天下至巧，亦猶須規矩也。}

「師曠之聰，不以六律，不能正五音。」\jiazhu[1]{師曠，晉平公之樂大師也。其聽至聰，不用六律，不能正五音。六律，陽律太蔟、姑洗、蕤賓、夷則、無射、黃鍾也。五音，宮、商、角、徵、羽也。}

「堯舜之道，不以仁政，不能平治天下。」\jiazhu[1]{當行仁恩之政，天下乃可平也。}

「今有仁心仁聞，而民不被其澤，不可法於後世者，不行先王之道也。」\jiazhu[1]{仁心，性仁也。仁聞，仁聲遠聞也。雖然，猶須行先王之道，使百姓被澤，乃可爲後法也。}

「故曰：徒善不足以爲政，徒法不能以自行。」\jiazhu[1]{但有善心而不行之，不足以爲政；但有善法度而不施之，法度亦不能獨自行也。}

「《詩》云：『不愆不忘，率由舊章。』遵先王之法而過者，未之有也。」\jiazhu[1]{《詩·大雅·嘉樂》之篇。愆，過也。所行不過差矣，不可忘者，以其循用舊故文章，遵用先王之法度，未聞有過也。}

「聖人既竭目力焉，繼之以規矩準繩，以爲方員平直，不可勝用也。」\jiazhu[1]{盡己目力，續以四者，方員平直可得而知審，故用之不可勝極也。}

「既竭耳力焉，繼之以六律，正五音，不可勝用也。」\jiazhu[1]{音須律而正也。}

「既竭心思焉，繼之以不忍人之政，而仁覆天下矣。」\jiazhu[1]{盡心欲行恩，繼以不忍加惡於人之政，則天下被覆衣之仁也。}

「故曰：爲高必因丘陵，爲下必因川澤。爲政不因先王之道，可謂智乎？」\jiazhu[1]{言因自然，則用力少而成功多矣。}

「是以惟仁者宜在高位。不仁而在高位，是播其惡於衆也。」\jiazhu[1]{仁者能由先王之道；不仁，逆道，則自播揚其惡於衆人也。}

「上無道揆也，下無法守也，朝不信道，工不信度，君子犯義，小人犯刑，國之所存者幸也。」\jiazhu[1]{言君無道術可以揆度天意，臣無法度可以守職奉命。朝廷之士不信道德，百工之作不信度量。君子觸義之所禁，謂學士當行君子之道也。小人觸刑，愚人罹於密網也。此亡國之政，然而國存者，僥倖耳，非其道也。}

「故曰：城郭不完，兵甲不多，非國之災也；田野不辟，貨財不聚，非國之害也。上無禮，下無學，賊民興，喪無日矣。」\jiazhu[1]{言君不知禮，臣不學法度，無以相檢制，則賊民興，亡在朝夕，無復有期日。言國無禮義必亡。}

「《詩》曰：『天之方蹶，無然泄泄。』泄泄猶沓沓也。事君無義，進退無禮，言則非先王之道者，猶沓沓也。」\jiazhu[1]{《詩·大雅·板》之篇。天，謂王者。蹶，動也。言天方動女，無敢沓沓但爲非義非禮，背棄先王之道而不相匡正也。}

「故曰：責難於君謂之恭，陳善閉邪謂之敬。吾君不能謂之賊。」\jiazhu[1]{人臣之道，當進君於善，責難爲之事，使君勉之，謂行堯舜之仁，是爲恭臣。陳善法以禁閉君之邪心，是爲敬君。言吾君不肖不能行善，因不諫正，此爲賊其君也。}\par \jiazhu[1]{章指言：雖有巧智，猶須法度。國由先王，禮義爲要。不仁在位，播越其惡。誣君不諫，故謂之賊。明上下相須，而道化行也。}

\section{論法堯舜}

孟子曰：「規矩，方員之至也；聖人，人倫之至也。」\jiazhu[1]{至，極也。人事之善者，莫大取法於聖人，猶方員須規矩也。}

「欲爲君，盡君道；欲爲臣，盡臣道。二者皆法堯舜而已矣。」\jiazhu[1]{堯舜之爲君臣，道備。}

「不以舜之所以事堯事君，不敬其君者也；不以堯之所以治民治民，賊其民者也。」\jiazhu[1]{言舜之事堯，敬之至也；堯之治民，愛之盡也。}

「孔子曰：『道二，仁與不仁而已矣。』暴其民甚，則身弑國亡；不甚，則身危國削。名之曰『幽厲』，雖孝子慈孫，百世不能改也。」\jiazhu[1]{仁則國安，不仁則國危亡。甚謂桀紂，不甚謂幽厲。厲王流于彘，幽王滅於戲，可謂身危國削矣。名之，謂謚之也。謚以幽厲，以章其惡，百世傳之，孝子慈孫何能改也？}

「《詩》云：『殷鑒不遠，在夏后之世。』此之謂也。」\jiazhu[1]{《詩·大雅·蕩》之篇也。殷之所鑒視，近在夏后之世耳。以前代善惡爲明鏡也，欲使周亦鑒於殷之所以亡也。}\par \jiazhu[1]{章指言：法則堯舜，以爲規矩；鑒戒桀紂，避遠危殆。名謚一定，千載而不可改也。}

\section{論仁與不仁}

孟子曰：「三代之得天下也以仁，其失天下也以不仁。國之所以廢興存亡者亦然。」\jiazhu[1]{三代，夏、殷、周。國，謂公侯之國。存亡在仁與不仁也。}

「天子不仁，不保四海；諸侯不仁，不保社稷；卿大夫不仁，不保宗廟；士庶人不仁，不保四體。今惡死亡而樂不仁，是由惡醉而強酒。」\jiazhu[1]{保，安也。四體，身之四肢。強酒則必醉也。}\par \jiazhu[1]{章指言：人所以安，莫若爲仁。惡而弗去，患必及身。自上達下，其道一焉。}

\section{論反求諸己}

孟子曰：「愛人不親，反其仁；治人不治，反其智；禮人不答，反其敬。行有不得者，皆反求諸己。其身正而天下歸之。」\jiazhu[1]{反其仁：己仁獨未至邪？反其智：己智獨未足邪？反其敬：己敬獨未恭邪？反求諸身，身己正，則天下歸就之，服其德也。}

「《詩》云：『永言配命，自求多福。』」\jiazhu[1]{此詩已見上篇，其義同。}\par \jiazhu[1]{章指言：行有不得於人，一求諸身，責己之道也。改行飭躬，福則至矣。}

\section{論天下國家之本}

孟子曰：「人有恒言，皆曰『天下國家』。」\jiazhu[1]{恒，常也。人之常語也。天下，謂天子之所主；國，謂諸侯之國；家，謂卿大夫之家也。}

「天下之本在國，國之本在家，家之本在身。」\jiazhu[1]{治天下者，不得良諸侯，無以爲本；治其國者，不得良卿大夫，無以爲本；治其家者，不得良身，無以爲本也。}\par \jiazhu[1]{章指言：天下國家，各依其本。本正則立，本傾則踣。雖曰常言，必須敬慎也。}

\section{論爲政得巨室}

孟子曰：「爲政不難，不得罪於巨室。」\jiazhu[1]{巨室，大家也，謂賢卿大夫之家，人所則效者。言不難者，但不使巨室罪之，則善也。}

「巨室之所慕，一國慕之；一國之所慕，天下慕之。故沛然德教溢乎四海。」\jiazhu[1]{慕，思也。賢卿大夫，一國思隨其所善惡。一國思其善政，則天下思以爲君矣。沛然大洽，德教可以滿溢於四海之內。}\par \jiazhu[1]{章指言：天下傾心，思慕嚮善。巨室不罪，咸以爲表。德之流行，可以充四海也。}

\section{論有道無道}

孟子曰：「天下有道，小德役大德，小賢役大賢；天下無道，小役大，弱役強。斯二者，天也。順天者存，逆天者亡。」\jiazhu[1]{有道之世，小德小賢樂爲大德大賢役，服於賢德也。無道之時，小國弱國畏懼而役於大國強國也。此二者，天時所遭也，當順從之，不當逆也。}

「齊景公曰：『既不能令，又不受命，是絕物也。』涕出而女於吳。」\jiazhu[1]{齊景公，齊侯，景，謚也。言諸侯既不能令告鄰國使之進退，又不能事大國往受教命，是所以自絕於物。物，事也。大國不與之通朝聘之事也。吳，蠻夷也，時爲強國，故齊侯畏而恥之，泣涕而與爲婚。}

「今也小國師大國而恥受命焉，是猶弟子而恥受命於先師也。」\jiazhu[1]{今小國以大國爲師，學法度焉，而恥受命教，不從其進退，譬猶弟子不從師也。}

「如恥之，莫若師文王。師文王，大國五年，小國七年，必爲政於天下矣。」\jiazhu[1]{文王行仁政，以移殷民之心，使皆就之。今師效文王，大國不過五年，小國七年，必得政於天下矣。文王時難，故百年乃洽；今之時易。文王由百里起，今大國乃踰千里，過之十倍有餘，故五年足以爲政；小國差之，故七年。}

「《詩》云：『商之孫子，其麗不億。上帝既命，侯于周服。侯服于周，天命靡常。殷士膚敏，祼將于京。』」\jiazhu[1]{《詩·大雅·文王》之篇。麗，億，數也。言殷帝之子孫，其數雖不伹億萬人，天既命之，惟服於周。殷之美士，執祼暢之禮，將事於京師，若微子者。膚，大；敏，達也。此天命之無常也。}

「孔子曰：『仁不可爲衆也。夫國君好仁，天下無敵。』」\jiazhu[1]{孔子云：行仁者，天下之衆不能當也。諸侯有好仁者，天下無敢與之爲敵。}

「今也欲無敵於天下而不以仁，是猶執熱而不以濯也。《詩》云：『誰能執熱，逝不以濯？』」\jiazhu[1]{《詩·大雅·桑柔》之篇。誰能持熱而不以水濯其手？喻爲國誰能違仁而無敵也。}\par \jiazhu[1]{章指言：遭衰逢亂，屈服強大。據國行仁，天下莫敵。雖有億衆，無德不親。執熱須濯，明不可違仁也。}

\section{論不仁者}

孟子曰：「不仁者，可與言哉？安其危而利其菑，樂其所以亡者。不仁而可與言，則何亡國敗家之有？」\jiazhu[1]{言不仁之人，以其所以爲危者，反以爲安；必以惡見亡，而樂行其惡。如使其能從諫從善，可與言議，則天下何有亡國敗家也？}

「有孺子歌曰：『滄浪之水清兮，可以濯我纓；滄浪之水濁兮，可以濯我足。』孔子曰：『小子聽之！清斯濯纓，濁斯濯足矣。自取之也。』」\jiazhu[1]{孺子，童子也。小子，孔子弟子也。清濯所用，尊卑若此，自取之，喻人善惡見尊賤乃如此。}

「夫人必自侮，然後人侮之；家必自毀，然後人毀之；國必自伐，然後人伐之。」\jiazhu[1]{人先自爲可侮慢之行，故見侮慢也；家先自爲可毀壞之道，故見毀也；國先自爲可誅伐之政，故見伐也。}

「《太甲》曰：『天作孽，猶可違；自作孽，不可活。』此之謂也。」\jiazhu[1]{已見上篇，說同。}\par \jiazhu[1]{章指言：人之安危，皆由於己。先自毀伐，人乃攻討，甚於天孽。敬慎而已，如臨深淵，戰戰恐慄也。}

\section{論得失天下}

孟子曰：「桀紂之失天下也，失其民也；失其民者，失其心也。」\jiazhu[1]{失其民之心，則天下畔之。簞食壺漿以迎武王之師是也。}

「得天下有道：得其民，斯得天下矣；得其民有道：得其心，斯得民矣；得其心有道：所欲與之聚之，所惡勿施爾也。」\jiazhu[1]{欲得民心，聚其所欲而與之。爾，近也。勿施行其所惡，使民近，則民心可得矣。}

「民之歸仁也，猶水之就下、獸之走壙也。故爲淵驅魚者，獺也；爲叢驅爵者，鸇也；爲湯武驅民者，桀與紂也。」\jiazhu[1]{民之思明君，猶水樂埤下，獸樂壙野。驅之則歸其所樂。獺，獱也。鸇，土鸇也。}

「今天下之君有好仁者，則諸侯皆爲之驅矣。雖欲無王，不可得已。」\jiazhu[1]{故云諸侯好爲仁者，驅民若此也。湯武行之矣，如有則之者，雖欲不王，不可得也。}

「今之欲王者，猶七年之病求三年之艾也。苟爲不畜，終身不得。苟不志於仁，終身憂辱，以陷於死亡。」\jiazhu[1]{今之諸侯欲行王道，而不積其德，如至七年病而却求三年時艾，當畜之乃可得，以三年時不畜藏之，至七年而欲卒求之，何可得乎？艾可以爲灸人病，乾久益善，故以爲喻。志仁者亦久行之，不行之則憂辱以陷死亡，桀紂是也。}

「《詩》云：『其何能淑，載胥及溺。』此之謂也。」\jiazhu[1]{《詩·大雅·桑柔》之篇。淑，善也。載，辭也。胥，相也。刺時君臣何能爲善乎？但相與爲沈溺之道也。}\par \jiazhu[1]{章指言：水性趨下，民樂歸仁。桀紂之驅，使就其君。三年之艾，畜而可得。一時欲仁，猶將沉溺，所以明鑒戒也。}

\section{論自暴自棄}

孟子曰：「自暴者，不可與有言也；自棄者，不可與有爲也。言非禮義，謂之自暴也；吾身不能居仁由義，謂之自棄也。」\jiazhu[1]{言人尚自暴自棄，何可與有言有爲？}

「仁，人之安宅也；義，人之正路也。曠安宅而弗居，舍正路而不由，哀哉！」\jiazhu[1]{曠，空；舍，縱；哀，傷也。弗由居是者，是可哀傷哉。}\par \jiazhu[1]{章指言：曠仁舍義，自暴棄之道也。}

\section{論道在邇而求諸遠}

孟子曰：「道在邇而求諸遠，事在易而求之難。人人親其親，長其長，而天下平。」\jiazhu[1]{邇，近也。道在近而患人求之遠也，事在易而苦人求之難也。謂不親其親，不事其長，故其事遠而難也。}\par \jiazhu[1]{章指言：親親敬長，近取諸己，則邇而易也。}

\section{論誠身之道}

孟子曰：「居下位而不獲於上，民不可得而治也。獲於上有道：不信於友，弗獲於上矣。信於友有道：事親弗悅，弗信於友矣。悅親有道：反身不誠，不悅於親矣。誠身有道：不明乎善，不誠其身矣。」\jiazhu[1]{言人求上之意，先從己始，本之於心。心不正而得人意者，未之有也。}

「是故誠者，天之道也；思誠者，人之道也。至誠而不動者，未之有也；不誠，未有能動者也。」\jiazhu[1]{授人誠善之性者，天也，故曰天道。思行其誠以奉天者，人道也。至誠則動金石，不誠則鳥獸不可親狎，故曰未有能動者也。}\par \jiazhu[1]{章指言：事上得君，乃可臨民；信友悅親，本在於身。是以曾子三省，《大雅》矜矜，以誠爲貴也。}

\section{論養老歸仁}

孟子曰：「伯夷辟紂，居北海之濱，聞文王作，興曰：『盍歸乎來！吾聞西伯善養老者。』」\jiazhu[1]{伯夷讓國，遭紂之世，辟之隱遁北海之濱。聞文王起興王道，盍歸乎來，歸周也。}

「太公辟紂，居東海之濱，聞文王作，興曰：『盍歸乎來！吾聞西伯善養老者。』」\jiazhu[1]{太公，呂望也，亦辟紂世，隱居東海，曰聞西伯養老。二人皆老矣，往歸文王也。}

「二老者，天下之大老也，而歸之，是天下之父歸之也。天下之父歸之，其子焉往？」\jiazhu[1]{此二老猶天下之父也，其餘皆天下之子耳。子當隨父，二父往矣，子將安如？言皆將往也。}

「諸侯有行文王之政者，七年之內，必爲政於天下矣。」\jiazhu[1]{今之諸侯如有能行文王之政者，七年之間，必足以爲政矣。天以七紀，故七年。文王時難，故久；衰周時易，故速也。上章言大國五年者，大國地廣人衆，易以行善，故五年足以治也。}\par \jiazhu[1]{章指言：養老尊賢，國之上務。文王勤之，二老遠至。父來子從，天之順道。七年爲政，以勉諸侯，欲使庶幾於行善也。}

\section{論不行仁政而富之}

孟子曰：「求也爲季氏宰，無能改於其德，而賦粟倍他日。孔子曰：『求非我徒也，小子鳴鼓而攻之可也。』」\jiazhu[1]{求，孔子弟子冉求。季氏，魯卿季康子。宰，家臣。小子，弟子也。孔子以冉求不能改季氏使從善，爲之多斂賦粟，故欲使弟子鳴鼓以聲其罪，而攻伐責讓之，曰求非我徒，疾之也。}

「由此觀之，君不行仁政而富之，皆棄於孔子者也。況於爲之強戰？爭地以戰，殺人盈野；爭城以戰，殺人盈城。此所謂率土地而食人肉，罪不容於死。」\jiazhu[1]{孔子棄富不仁之君者，況於爭地而殺人滿之乎？此若率土地使食人肉也，言其罪大，死刑不足以容之。}

「故善戰者服上刑，連諸侯者次之，辟草萊、任土地者次之。」\jiazhu[1]{孟子言天道重生，戰者殺人，故使善戰者服上刑。上刑，重刑也。連諸侯，合從者也，罪次善戰者。辟草任地，不務修德而富國者，罪次合從連橫之人也。}\par \jiazhu[1]{章指言：聚斂富君，棄於孔子，冉求行之。同聞鳴鼓，以戰殺民，土食人肉，罪不容死，以爲大戮，重人命之至也。}

\section{論觀人之眸子}

孟子曰：「存乎人者，莫良於眸子。眸子不能掩其惡。」\jiazhu[1]{眸子，目瞳子也。存人，存在人之善惡也。}

「胸中正，則眸子瞭焉；胸中不正，則眸子眊焉。」\jiazhu[1]{瞭，明也。眊者，蒙蒙目不明之貌。}

「聽其言也，觀其眸子，人焉廋哉？」\jiazhu[1]{廋，匿也。聽言察目，言正視端，人情可見，安可匿哉？}\par \jiazhu[1]{章指言：目可神候，精之所在，存而察之，善惡不隠，知人之道，斯爲審矣。}

\section{論恭儉}

孟子曰：「恭者不侮人，儉者不奪人。侮奪人之君，惟恐不順焉，惡得爲恭儉？」\jiazhu[1]{爲恭敬者不侮慢人，爲廉儉者不奪取人。有好侮奪人之君，有貪陵之性，恐人不順從其所欲，安得爲恭儉之行也？}

「恭儉豈可以聲音笑貌爲哉？」\jiazhu[1]{恭儉之人，儼然無欲，自取其名，豈可以和聲諂笑之貌強爲之哉？}\par \jiazhu[1]{章指言：人君恭儉，率下移風；人臣恭儉，明其廉忠。侮奪之惡，何由干之而錯其心？}

\section{淳于髡論權}

淳于髡曰：「男女授受不親，禮與？」\jiazhu[1]{淳于髡，齊人也。問禮男女不相親授。}

孟子曰：「禮也。」\jiazhu[1]{禮不親授。}

曰：「嫂溺，則援之以手乎？」\jiazhu[1]{髡曰：見嫂溺水，則當以手牽援之否邪？}

曰：「嫂溺不援，是豺狼也。」\jiazhu[1]{孟子曰：人見嫂溺不援出，是爲豺狼之心也。}

「男女授受不親，禮也；嫂溺，援之以手者，權也。」\jiazhu[1]{孟子告髡曰：此權也。權者，反經而善也。}

曰：「今天下溺矣，夫子之不援，何也？」\jiazhu[1]{髡曰：今天下之道溺矣，夫子何不援之？}

曰：「天下溺，援之以道；嫂溺，援之以手。子欲手援天下乎？」\jiazhu[1]{孟子曰：當以道援天下，而道不得行，子欲使我以手援天下乎？}\par \jiazhu[1]{章指言：權時之義，嫂溺援手；君子大行，拯世以道，道之指也。}

\section{論君子不教子}

公孫丑曰：「君子之不教子，何也？」\jiazhu[1]{問父子不親教何也？}

孟子曰：「勢不行也。教者必以正，以正不行，繼之以怒。繼之以怒，則反夷矣。『夫子教我以正，夫子未出於正也。』則是父子相夷也。父子相夷，則惡矣。」\jiazhu[1]{父親教子，其勢不行。教以正道而不能行，則責怒之。夷，傷也，父子相責怒則傷義矣。一說曰：父子反自相非若夷狄也。子之心責其父云：夫子教我以正道，而夫子之身未必自行正道也。執此意，則爲反夷矣，故曰惡也。}

「古者易子而教之，父子之間不責善。責善則離，離則不祥莫大焉。」\jiazhu[1]{易子而教，不欲自相責以善也。父子主恩，離則不祥莫大焉。}\par \jiazhu[1]{章指言：父子至親，相責離恩；易子而教，相成以仁。教之義也。}

\section{論事親守身}

孟子曰：「事孰爲大？事親爲大。守孰爲大？守身爲大。不失其身而能事其親者，吾聞之矣；失其身而能事其親者，吾未之聞也。」\jiazhu[1]{事親，養親也。守身，使不陷於不義也。失仁義，則何能事父母乎？}

「孰不爲事？事親，事之本也。孰不爲守？守身，守之本也。」\jiazhu[1]{先本後末，事守乃立也。}

「曾子養曾晳，必有酒肉。將徹，必請所與；問有餘，必曰『有』。曾晳死，曾元養曾子，必有酒肉。將徹，不請所與；問有餘，曰『亡矣』。將以復進也。此所謂養口體者也。若曾子，則可謂養志也。事親若曾子者，可也。」\jiazhu[1]{將徹，請所與，問曾晳所欲與子孫所愛者也。必曰有，恐違親意也，故曰養志。曾元曰無，欲以復進曾子也，不求親意，故曰養口體也。事親之道，當如曾子之法，乃爲至孝也。}\par \jiazhu[1]{章指言：上孝養志，下孝養體。曾參事親，可謂至矣。孟子言之，欲令後人則曾子也。}

\section{論格君心之非}

孟子曰：「人不足與適也，政不足閒也。惟大人爲能格君心之非。」\jiazhu[1]{適，過也。《詩》云：「室人交徧謫我。」閒，非。格，正也。時皆小人居位，不足過責也；政教不足復非訧；獨得大人爲輔臣，乃能正君之非法度也。}

「君仁，莫不仁；君義，莫不義；君正，莫不正。一正君而國定矣。」\jiazhu[1]{正君之身，一國定矣。欲使大人正之。}\par \jiazhu[1]{章指言：小人爲政，不足閒非；賢臣正君，使握道機。君正國定，下不邪侈，將何閒也？}

\section{論不虞之譽與求全之毀}

孟子曰：「有不虞之譽，有求全之毀。」\jiazhu[1]{虞，度也。言人之行，有不度其將有名譽而得者，若尾生本與婦人期於梁下，不度水之卒至，遂至沒溺而獲守信之譽。求全之毀，若陳不瞻將赴君難，聞金鼓之聲，失氣而死，可謂欲求全其節，而反有怯弱之毀者也。}\par \jiazhu[1]{章指言：不虞獲譽，不可爲戒；求全受毀，未足懲咎。君子正行，不由斯二者也。}

\section{論易其言}

孟子曰：「人之易其言也，無責耳矣。」\jiazhu[1]{人之輕易其言，不得失言之咎責也。一說：人之輕易不肯諫正君者，以其不在言責之位者也。}\par \jiazhu[1]{章指言：言出於身，駟不及舌；不惟其責，則易之矣。}

\section{論好爲人師}

孟子曰：「人之患在好爲人師。」\jiazhu[1]{人之所患，患於不知己未有可師，而好爲人師者，惑也。}\par \jiazhu[1]{章指言：君子好謀而成，臨事而懼，時然後言，畏失言也。故曰師哉師哉，桐子之命。不慎則有患矣。}

\section{樂正子見孟子}

樂正子從於子敖之齊。樂正子見孟子。\jiazhu[1]{魯人樂正克，孟子弟子也。從於齊之右師子敖。子敖使而之魯，樂正子隨之來之齊也。孟子在齊，樂正子見之也。}

孟子曰：「子亦來見我乎？」\jiazhu[1]{孟子見其來見遲，故云亦來也。}

曰：「先生何爲出此言也？」\jiazhu[1]{樂正子曰：先生何爲非克而出此言？}

曰：「子來幾日矣？」\jiazhu[1]{孟子問子來幾日乎？}

曰：「昔者。」\jiazhu[1]{克曰：昔者來至。昔者，往也，謂數日之間也。}

曰：「昔者，則我出此言也，不亦宜乎？」\jiazhu[1]{孟子曰：昔者來至，而今乃來，我出此言，亦其宜也。孟子重愛樂正子，欲亟見之，思深望重也。}

曰：「舍館未定。」\jiazhu[1]{克曰：所止舍館未定，故不即來也。館，客舍。}

曰：「子聞之也，舍館定，然後求見長者乎？」\jiazhu[1]{孟子曰：子聞見長者之禮，當須舍館定乃見之乎？}

曰：「克有罪。」\jiazhu[1]{樂正子謝過服罪也。}\par \jiazhu[1]{章指言：尊重道，敬賢事長，人之大綱。樂正子好善，故孟子譏之，責賢者備也。}

\section{孟子謂樂正子}

孟子謂樂正子曰：「子之從於子敖來，徒餔啜也。我不意子學古之道，而以餔啜也。」\jiazhu[1]{子敖，齊之貴人右師王驩也。學而不行其道，徒食飲而已，謂之餔啜也。樂正子本學古聖人之道，而今隨從貴人，無所匡正，故言不意子但餔啜也。}\par \jiazhu[1]{章指言：學優則仕，仕以行道；否則隱逸，免置窮處。餔啜沈浮，君子不與。是以孟子咨嗟樂正子也。}

\section{論不孝有三}

孟子曰：「不孝有三，無後爲大。」\jiazhu[1]{於禮有不孝者三事：謂阿意曲從，陷親不義，一不孝也；家貧親老，不爲禄仕，二不孝也；不娶無子，絕先祖祀，三不孝也。三者之中，無後爲大。}

「舜不告而娶，爲無後也。君子以爲猶告也。」\jiazhu[1]{舜懼無後，故不告而娶。君子知舜告焉不得而娶，娶而告父母，禮也；舜不以告，權也，故曰猶告，與告同也。}\par \jiazhu[1]{章指言：量其輕重，無後不可。是以大舜受堯二女。夫三不孝，蔽者所闇，至於大聖，卓然匪疑，所以垂法也。}

\section{論仁義禮樂之實}

孟子曰：「仁之實，事親是也；義之實，從兄是也；智之實，知斯二者弗去是也。」\jiazhu[1]{事皆有實。事親從兄，仁義之實也。知仁義所用而不去之，則智之實也。}

「禮之實，節文斯二者是也；樂之實，樂斯二者。」\jiazhu[1]{禮樂之實，節文事親從兄，使不失其節，而文其禮敬之容，而中心樂之也。}

「樂則生矣，生則惡可已也。惡可已也，則不知足之蹈之、手之舞之。」\jiazhu[1]{樂此事親從兄出於中心，則樂生其中矣。樂生之至，安可已也？豈從自覺足蹈節、手舞曲哉？}\par \jiazhu[1]{章指言：仁義之本，在於孝弟。孝弟之至，通於神明，況於歌舞？不能自知，蓋有諸中，形於外也。}

\section{論大孝}

孟子曰：「天下大悅而將歸己，視天下悅而歸己，猶草芥也，惟舜爲然。」\jiazhu[1]{舜不以天下將歸己爲樂，號泣于天。}

「不得乎親，不可以爲人；不順乎親，不可以爲子。舜盡事親之道而瞽瞍厎豫，瞽瞍厎豫而天下化，瞽瞍厎豫而天下之爲父子者定。此之謂大孝。」\jiazhu[1]{舜以不順親意爲非人子。厎，致也；豫，樂也。瞽瞍，頑父也。盡其孝道而頑父致樂，使天下化之爲父子之道者定也。}\par \jiazhu[1]{章指言：以天下之貴富爲不若得意於親，故能懷協頑嚚，厎豫而欣，天下化之，父子加親。故稱盛德者必百世祀，無與比崇也。}
